\chapter{Understanding MQTT and Sparkplug~B}
\label{ch:understanding}

This chapter explains \emph{how Sparkplug~B works}, for a reader who has not necessarily
worked with MQTT before. It is background, not product documentation: it describes the
protocol, so that the rest of this manual --- and the topics you will see on your broker
--- read as a coherent design rather than as a list of conventions.

\begin{note}
\textbf{This chapter describes the specification. Chapter~\ref{ch:contract} describes
\prog.} The two are deliberately kept apart, because conflating them is how a manual comes
to describe capabilities a product does not have. Wherever this chapter mentions \prog it
says so explicitly, and section~\ref{sec:sp-where-we-sit} at the end maps every mechanism
onto what is implemented, what is not, and where \prog knowingly departs from the norm.

The authority for \prog's behaviour is chapter~\ref{ch:contract} and the clause-by-clause
audit in \code{docs/sparkplug-conformance.md}; the authority for the protocol is the
specification itself, vendored in this repository (see
section~\ref{sec:sp-reading-the-norm}).
\end{note}

\section{MQTT, in the terms Sparkplug needs}

Sparkplug is not a replacement for MQTT --- it is a set of conventions \emph{on top of}
MQTT. Almost every Sparkplug mechanism is built out of one of five MQTT features, so those
five are worth understanding first. Everything else about MQTT can wait.

\subsection{Publish/subscribe, and the broker in the middle}

MQTT clients never talk to each other. They connect to a \textbf{broker} (the
specification calls it an \emph{MQTT Server}) and either \textbf{publish} messages to a
\emph{topic}, or \textbf{subscribe} to topics and receive whatever anyone publishes there.

\begin{center}
\begin{tikzpicture}[node distance=9mm]
  \node[actor] (edge) {Publisher\\ \scriptsize(e.g.\ \prog)};
  \node[brokeract, right=22mm of edge] (broker) {MQTT broker};
  \node[actor, right=22mm of broker, yshift=9mm]  (h1) {Subscriber\\ \scriptsize(Ignition)};
  \node[actor, right=22mm of broker, yshift=-11mm] (h2) {Subscriber\\ \scriptsize(anything else)};
  \draw[msg] (edge) -- node[above, font=\sffamily\scriptsize] {publish} (broker);
  \draw[msg] (broker) -- node[above, font=\sffamily\scriptsize, pos=0.55] {} (h1);
  \draw[msg] (broker) -- node[below, font=\sffamily\scriptsize, pos=0.55] {} (h2);
  \node[annot, below=7mm of broker, align=center]
       {the publisher does not know who is listening,\\ or whether anyone is at all};
\end{tikzpicture}
\end{center}

That decoupling is the point, and it is also the difficulty Sparkplug exists to solve. A
publisher gets no answer, so it cannot tell a healthy silence (nothing has changed) from a
broken one (it is disconnected, or the consumer is gone). There is no request/response and
no polling.

\subsection{Topics and wildcards}

A topic is a UTF-8 string cut into levels by forward slashes, e.g.\
\code{spBv1.0/Site/NDATA/Bridge}. Subscribers may use two wildcards:

\begin{center}
\begin{tabular}{@{}>{\ttfamily}c l p{7.6cm}@{}}
\toprule
\normalfont\textbf{Wildcard} & \textbf{Matches} & \textbf{Example} \\
\midrule
+ & exactly one level  & \code{spBv1.0/+/NDATA/+} --- every node's data, any group \\
\# & all remaining levels & \code{spBv1.0/\#} --- the whole Sparkplug namespace \\
\bottomrule
\end{tabular}
\end{center}

This is why Sparkplug forbids the characters \code{+}, \code{/} and \code{\#} inside any
identifier you choose (group, node or device): they would either split a level in two or
turn a name into a wildcard. Topics are also \textbf{case-sensitive} --- \code{a/b} and
\code{A/b} are different topics --- and Sparkplug inherits that, for metric names as well.
The specification recommends against identifiers that differ only in case, because many
host applications and most databases cannot tell them apart.

\subsection{Quality of Service: 0, 1 and 2}

Every publish carries a QoS level, which is a promise about \emph{delivery attempts}, not
about correctness:

\begin{center}
\begin{tabular}{@{}c l p{8.2cm}@{}}
\toprule
\textbf{QoS} & \textbf{Promise} & \textbf{Cost and consequence} \\
\midrule
0 & at most once  & Fire and forget. No acknowledgement, no retry, no queueing. The
                    publisher learns only that the message was handed to the transport. \\
1 & at least once & Acknowledged and retried until confirmed, so a message may arrive
                    \emph{twice}. The receiver must tolerate duplicates. \\
2 & exactly once  & A four-packet handshake. Slowest, and rarely used in Sparkplug. \\
\bottomrule
\end{tabular}
\end{center}

Sparkplug requires \textbf{QoS~0} for the data-carrying messages (birth and data), on the
grounds that a stale value retried later is worth less than a fresh one, and \textbf{QoS~1}
for the death certificate registered with the broker --- the one message whose loss is
never acceptable.

\begin{warning}
\textbf{At QoS~0 an application cannot honestly claim a message was \emph{delivered}.}
There is no acknowledgement to observe. It can only claim the transport accepted it for
transmission. That distinction is not pedantry --- it is the difference between a screen
that reports what happened and one that reports what was hoped for, and it is why
\prog's own delivery wording is worded the way it is (chapter~\ref{ch:contract}).
\end{warning}

\subsection{Retained messages --- the broker's memory}

A published message normally reaches only clients subscribed \emph{at that moment}. Set
the \textbf{retain} flag and the broker stores it instead: the last retained message on a
topic is delivered immediately to every future subscriber, for as long as the broker
lives.

\begin{warning}
\textbf{A retained message outlives its publisher, indefinitely.} It says what was last
published on a topic, not what is true now, and not that the publisher still exists. A
client that connected once and died years ago leaves its last retained message behind
forever.

This project has already paid for that distinction. An observation of the production
broker found four host identifiers with retained state; three of them turned out to be
residue from long-dead clients, and a careful reading of the retained set alone produced a
confident conclusion that a single restart then destroyed. \textbf{A snapshot of retained
messages describes the broker's memory, not any participant's behaviour} --- the full
record is in \code{docs/primary-host-state-observation.md}.
\end{warning}

Sparkplug uses retain for exactly one thing --- the host application's \code{STATE}
(section~\ref{sec:sp-primary-host}) --- and forbids it everywhere else, because a retained
measurement is a value replayed to a future subscriber with no way to judge its age.

\subsection{The Last Will and Testament}

This is the feature Sparkplug is built around. When a client connects it may hand the
broker a complete message --- topic, payload, QoS, retain flag --- inside the CONNECT
packet itself, with the instruction: \emph{if I disappear without saying goodbye, publish
this on my behalf.}

The will fires when the broker decides the client is gone: the socket closed, or the
keep-alive elapsed with no traffic. It does \textbf{not} fire if the client sends a clean
MQTT DISCONNECT packet --- that tells the broker the departure was intentional, and the
broker discards the will.

\begin{center}
\begin{tikzpicture}
  \node[actor]     (n) at (0,0)    {Edge Node};
  \node[brokeract] (b) at (5.0,0)  {Broker};
  \node[actor]     (h) at (10.0,0) {Host};
  \draw[lifeline] (n) -- (0,-3.1);
  \draw[lifeline] (b) -- (5.0,-3.1);
  \draw[lifeline] (h) -- (10.0,-3.1);

  \draw[msg]     (0,-0.9) -- node[mlbl]{CONNECT (will = NDEATH)} (5.0,-0.9);
  \node[annot]   at (2.5,-1.6) {\dots\ the node vanishes without a DISCONNECT \dots};
  \draw[msgdead] (5.0,-2.3) -- node[mlbl]{NDEATH} (10.0,-2.3);
  \node[annot, align=center] at (5.0,-2.95)
       {the broker publishes the certificate the node registered\\ before it had anything to report};
\end{tikzpicture}
\end{center}

\subsection{Keep-alive, and the half-open connection}

The client and broker agree a \textbf{keep-alive} interval at CONNECT. If the broker sees
no traffic for one and a half times that interval, it declares the client dead and fires
the will.

This bounds --- but does not eliminate --- the time a dead node can be displayed as live.
A socket closed cleanly is noticed at once. A \emph{half-open} connection is not: a
container network torn down without a FIN, or a firewall silently dropping a flow, leaves
the broker holding a connection that no longer exists until the keep-alive elapses. With a
30-second keep-alive that is up to 45 seconds of a dead node shown as healthy. It is the
reason a well-behaved node publishes its \emph{own} death certificate before a planned
stop instead of relying on the will alone.

\subsection{Clean session}

MQTT can remember a client between connections --- its subscriptions, and messages queued
while it was away. Sparkplug requires the opposite: an edge node MUST connect with
\textbf{clean session true} (MQTT~3.1.1), or clean start with a session expiry of zero
(MQTT~5.0). Session state is Sparkplug's own job, carried explicitly in birth and death
certificates, and a broker replaying a backlog of stale measurements after a reconnect
would undermine it.

\section{The problem Sparkplug solves}

Plain MQTT is stateless from the consumer's point of view. A SCADA system needs the
opposite: for every point on screen it must know not only the last value but
\textbf{whether that value is still current}. Legacy SCADA protocols get this from polling
--- \emph{I ask, I get an answer, therefore I know the state; and I must ask again,
because it may have changed.} Polling is exactly what MQTT was designed to avoid.

Sparkplug's answer has three parts, and they only work together:

\begin{enumerate}
  \item \textbf{Session state made explicit.} Every node announces itself with a
        \emph{birth certificate} and is guaranteed to be announced dead by a \emph{death
        certificate} --- registered as the MQTT will, so it arrives even if the node
        vanishes without warning.
  \item \textbf{Report by exception.} Because the consumer can trust that the session is
        alive, the node need only publish values \emph{when they change}. Silence becomes
        meaningful: it means nothing has changed, not that something is broken.
  \item \textbf{Self-describing payloads.} The birth carries the complete set of metrics
        with names, datatypes and current values, so a consumer needs no prior
        configuration to interpret what follows.
\end{enumerate}

\begin{note}
Point~2 is why the specification says data \emph{should not} be published periodically
(\code{tck-id-principles-rbe-recommended}). \prog does publish periodically today, and
records it as a deviation --- see section~\ref{sec:sp-where-we-sit}.
\end{note}

\section{The pieces of a Sparkplug infrastructure}

\begin{center}
\begin{tabular}{@{}l p{9.4cm}@{}}
\toprule
\textbf{Component} & \textbf{What it is} \\
\midrule
MQTT Server & An ordinary MQTT broker meeting a modest subset of the MQTT
              specification: QoS~0 and~1, full will-message support, full retain
              support. \\
Edge Node   & An MQTT client that manages a Sparkplug session and speaks for itself and
              for the devices beneath it. \textbf{This is \prog's role.} \\
Device      & A physical or logical thing behind an edge node --- a PLC, a sensor, a
              meter. It has no MQTT connection of its own; the node publishes on its
              behalf. \\
Host Application & Any client consuming Sparkplug messages. It must publish
              \code{STATE} messages announcing whether it is online. \\
Primary Host Application & A host application that a \emph{given edge node} has been
              configured to treat as decisive: the node changes its own behaviour
              depending on whether that host is online. The relationship is declared at
              the node, not at the host. \\
Sparkplug Aware MQTT Server & A broker that additionally stores every NBIRTH and DBIRTH
              passing through it and republishes them as retained messages on
              \code{\$sparkplug} topics, so a consumer arriving late can still collect
              them. A distinct and stronger conformance profile. \\
\bottomrule
\end{tabular}
\end{center}

\begin{note}
In this installation \prog is the Edge Node, each smart-me meter is a Device, Ignition
with MQTT~Engine is the Host Application, and the broker is Mosquitto --- which is
\emph{Sparkplug Compliant} but \textbf{not Sparkplug Aware}. That last point has a real
consequence, drawn out in section~\ref{sec:sp-where-we-sit}.
\end{note}

\section{The topic namespace}

Every Sparkplug topic has the same shape:

\begin{center}
\code{namespace/group\_id/message\_type/edge\_node\_id/[device\_id]}
\end{center}

\begin{itemize}
  \item \textbf{namespace} is the fixed string \code{spBv1.0}. It declares both the topic
        grammar and the payload encoding.
  \item \textbf{group\_id} logically groups edge nodes --- a site, a production line, a
        pipeline segment.
  \item \textbf{message\_type} is one of the nine verbs below.
  \item \textbf{edge\_node\_id} identifies the node. The pair
        \code{group\_id/edge\_node\_id} must be unique across the whole infrastructure.
  \item \textbf{device\_id} identifies a device beneath the node. It is present on exactly
        four of the nine verbs and forbidden on the rest.
\end{itemize}

Group, node and device identifiers may be any UTF-8 string \emph{except} that they may not
contain \code{+}, \code{/} or \code{\#}. Device identifiers must be unique within their
node, but may repeat across different nodes.

\subsection{The nine verbs}

\begin{center}
\begin{tabular}{@{}>{\ttfamily}l c l p{6.3cm}@{}}
\toprule
\normalfont\textbf{Verb} & \textbf{Device?} & \textbf{Direction} & \textbf{Meaning} \\
\midrule
NBIRTH & no  & node $\rightarrow$ host & Node birth certificate. Declares every metric the
                                         node owns, with datatypes and current values. \\
NDEATH & no  & node $\rightarrow$ host & Node death certificate. Registered as the MQTT
                                         will; may also be published explicitly. \\
DBIRTH & yes & node $\rightarrow$ host & Device birth certificate, one per device. \\
DDEATH & yes & node $\rightarrow$ host & A device has stopped reporting, while the node
                                         itself is still alive. \\
NDATA  & no  & node $\rightarrow$ host & Node-level measurements. \\
DDATA  & yes & node $\rightarrow$ host & Device-level measurements. \\
NCMD   & no  & host $\rightarrow$ node & Command to the node. Includes the mandatory
                                         rebirth request. \\
DCMD   & yes & host $\rightarrow$ node & Command to a device --- typically writing an
                                         output. \\
STATE  & \multicolumn{2}{l}{\emph{host, own topic}} & The host application's own
                                         online/offline announcement. \\
\bottomrule
\end{tabular}
\end{center}

\code{STATE} is the exception to every rule here: its topic is
\code{spBv1.0/STATE/sparkplug\_host\_id} --- no group, no verb position in the usual sense
--- and its payload is not a Sparkplug payload at all. It is covered in
section~\ref{sec:sp-primary-host}.

\section{The payload}

With the single exception of \code{STATE}, every Sparkplug~B payload is a
\textbf{Google Protocol Buffers} message: a compact binary encoding, not human-readable on
the wire. You cannot read it with \code{mosquitto\_sub} alone; you need a decoder.

\subsection{The envelope}

\begin{center}
\begin{tabular}{@{}>{\ttfamily}l l p{7.9cm}@{}}
\toprule
\normalfont\textbf{Field} & \textbf{Type} & \textbf{Meaning} \\
\midrule
timestamp & uint64 & Milliseconds since the Unix epoch, \textbf{UTC}, for the message
                     as a whole. \\
seq       & uint64 & Sequence number, 0--255, used by the consumer to detect a gap.
                     See below. \\
metrics   & array  & The measurements and declarations --- the substance. \\
uuid      & string & Optional; identifies a schema or encoding for the \code{body}. \\
body      & bytes  & Optional; arbitrary binary data, rarely used. \\
\bottomrule
\end{tabular}
\end{center}

\subsection{The sequence number, and what it is for}

Every message from an edge node carries \code{seq} --- \emph{except} NDEATH, which
deliberately does not, because the broker may publish it long after the node stopped
counting.

The NBIRTH sets a starting value between 0 and 255. Every subsequent message increments it
by one, wrapping from 255 back to 0. A consumer that receives 41, 42, 44 knows message 43
existed and did not arrive, and can ask for a rebirth rather than silently carrying a hole
in its data. \textbf{That is the whole purpose:} the number is not a timestamp and not an
identifier, it is a gap detector.

\subsection{bdSeq --- pairing a death to its birth}

\code{bdSeq} (``birth/death sequence'') is a metric, not an envelope field, and it appears
in exactly two places: the birth certificate and the death certificate of the same
session. They carry the \emph{same} value; the next MQTT CONNECT uses the next value,
wrapping 255 to 0.

It exists because MQTT gives no ordering guarantee \emph{between} sessions. A death
certificate from a session that ended can be delivered after the next session's birth ---
and a consumer that acted on it would mark a healthy node dead. Matching \code{bdSeq}
makes that impossible to get wrong: a death whose number does not match the current birth
is stale and must be ignored.

\begin{note}
Because of this, a consumer must \textbf{not} use the NDEATH's timestamp to pair it with a
birth. The specification says so explicitly: the will's payload was built before the
CONNECT, so its timestamp is the connect time, not the death time. \code{bdSeq} is the
only correct pairing key.
\end{note}

\subsection{Metrics}

A metric is Sparkplug's equivalent of a SCADA tag:

\begin{center}
\begin{tabular}{@{}>{\ttfamily}l p{10.0cm}@{}}
\toprule
\normalfont\textbf{Field} & \textbf{Meaning} \\
\midrule
name       & A slash-delimited UTF-8 name, where slashes denote folders:
             \code{Node Control/Rebirth} is metric \code{Rebirth} in folder
             \code{Node Control}. Required, unless aliases are in use. \\
alias      & An optional unsigned integer standing in for the name. Declared alongside
             the name in a birth, then used \emph{instead of} the name in data messages
             to save bandwidth. \\
timestamp  & When this metric's value was \textbf{captured} --- not when it was
             published. UTC, milliseconds. \\
datatype   & One of the enumerated Sparkplug types. Required in births, and deliberately
             omitted from data messages, because the birth already established it. \\
is\_null      & Explicitly marks the value as null, rather than omitting it ambiguously. \\
is\_historical & Marks a value as \emph{not} a current reading --- replayed from a store,
             or backfilled after an outage. A consumer must not treat it as live. \\
is\_transient & Of interest on screen, but should not be written to a historian. \\
properties & A set of custom key/value pairs attached to the metric. \textbf{This is where
             quality and engineering units travel.} \\
metadata   & Descriptive information for complex or binary types. \\
value      & The value itself, in one of the types below. \\
\bottomrule
\end{tabular}
\end{center}

\subsection{Datatypes}

Sparkplug enumerates its types by number, and the number --- not the name --- is what
travels on the wire.

\begin{center}
\begin{tabular}{@{}l p{9.3cm}@{}}
\toprule
\textbf{Group} & \textbf{Types (with their wire codes)} \\
\midrule
Integers & \code{Int8}~1, \code{Int16}~2, \code{Int32}~3, \code{Int64}~4,
           \code{UInt8}~5, \code{UInt16}~6, \code{UInt32}~7, \code{UInt64}~8 \\
Reals    & \code{Float}~9, \code{Double}~10 \\
Other basics & \code{Boolean}~11, \code{String}~12, \code{DateTime}~13, \code{Text}~14 \\
Structured & \code{UUID}~15, \code{DataSet}~16, \code{Bytes}~17, \code{File}~18,
             \code{Template}~19 \\
Property-only & \code{PropertySet}~20, \code{PropertySetList}~21 \\
Arrays   & \code{Int8Array}~22 \dots \code{DateTimeArray}~34 \\
\midrule
\code{Unknown}~0 & A placeholder for future expansion --- never a valid published type. \\
\bottomrule
\end{tabular}
\end{center}

Two of the structured types are worth naming even though \prog uses neither:
a \textbf{DataSet} carries a table --- named, typed columns and rows --- in a single
metric, and a \textbf{Template} defines a reusable composite type (a ``UDT''), so that a
pump with six attributes can be declared once and instantiated many times.

\section{The session, as a state machine}
\label{sec:sp-session}

This is the heart of Sparkplug. The order of operations is not a style preference --- each
step exists because a failure at that exact moment would otherwise go unreported.

\subsection{Establishing a session}

\begin{enumerate}
  \item \textbf{Choose the session number.} Take the next \code{bdSeq}.
  \item \textbf{Build the death certificate first.} Serialise the NDEATH for the session
        that has not yet started.
  \item \textbf{Register it as the will, inside the CONNECT packet.} Not after connecting
        --- inside. There must be no instant in which the node is connected and the broker
        holds no certificate for it.
  \item \textbf{Connect} (clean session, QoS~1 will, retain false on the will).
  \item \textbf{Publish NBIRTH} --- the first message, always. It declares every node-level
        metric with its datatype and current value, carries \code{seq} and the session's
        \code{bdSeq}, and must include the \code{Node Control/Rebirth} command metric.
  \item \textbf{Publish one DBIRTH per device}, each declaring that device's metrics.
  \item \textbf{Publish data} as values change, incrementing \code{seq} each time.
\end{enumerate}

\begin{center}
\begin{tikzpicture}[node distance=0mm]
  % lifelines
  \node[actor] (n) at (0,0) {Edge Node};
  \node[brokeract] (b) at (5.0,0) {Broker};
  \node[actor] (h) at (10.0,0) {Host};
  \draw[lifeline] (n) -- (0,-6.6);
  \draw[lifeline] (b) -- (5.0,-6.6);
  \draw[lifeline] (h) -- (10.0,-6.6);

  \draw[msg] (0,-0.9) -- node[mlbl]{CONNECT (will = NDEATH, QoS 1)} (5.0,-0.9);
  \draw[msg] (0,-1.7) -- (5.0,-1.7) node[midway, mlbl]{NBIRTH (seq 0, bdSeq n)};
  \draw[msg] (5.0,-1.7) -- (10.0,-1.7);
  \draw[msg] (0,-2.5) -- (5.0,-2.5) node[midway, mlbl]{DBIRTH (seq 1)};
  \draw[msg] (5.0,-2.5) -- (10.0,-2.5);
  \draw[msg] (0,-3.3) -- (5.0,-3.3) node[midway, mlbl]{DDATA (seq 2)};
  \draw[msg] (5.0,-3.3) -- (10.0,-3.3);
  \node[annot] at (5.0,-3.9) {\dots\ values only when they change \dots};

  \draw[msgdead] (0,-4.7) -- (5.0,-4.7) node[midway, mlbl]{NDEATH (explicit, planned stop)};
  \draw[msgdead] (5.0,-4.7) -- (10.0,-4.7);
  \node[annot, align=left] at (2.5,-5.35) {socket drops --- no DISCONNECT sent};
  \draw[msgdead] (5.0,-5.9) -- (10.0,-5.9) node[midway, mlbl]{NDEATH (the will, from the broker)};
  \node[annot, align=left] at (7.5,-6.45) {same bdSeq --- a second death is idempotent};
\end{tikzpicture}
\end{center}

\subsection{Ending a session}

There are two ways a session ends, and a well-behaved node arranges for both to be
visible:

\begin{itemize}
  \item \textbf{A planned stop.} The node publishes its own NDEATH, so the host is told
        immediately instead of waiting out the keep-alive. The specification says a node
        \emph{should} do this.
  \item \textbf{A hard death} --- crash, power loss, severed network. Nothing runs on the
        node, and the broker publishes the registered will.
\end{itemize}

The subtlety is that these are not alternatives. If a node publishes an explicit death and
then sends a clean MQTT DISCONNECT, the broker \textbf{discards the will} --- and any
future failure of that same connection goes unannounced. A node that wants both
protections publishes its death and then simply drops the connection, which produces
\textbf{two} death certificates carrying the same \code{bdSeq}. A consumer that has already
marked the node down ignores the second; there is no harm in it.

\section{Report by exception}

Because session state is explicit and guaranteed, a node does not need to repeat itself.
The births carry every current value; afterwards, only changes are published. The
specification puts it as a \emph{should not}: data should not be published periodically.

This matters more than it first appears. Under report-by-exception, an unchanged value and
a broken publisher look different --- the first produces silence with a live session, the
second produces a death certificate. Under periodic publishing they look the same, and the
consumer loses the ability to tell them apart from the data stream alone.

\section{Commands, and the rebirth handshake}
\label{sec:sp-rebirth}

Commands travel the other way: a host publishes \code{NCMD} to a node, or \code{DCMD} to a
device, typically to write an output.

One command is mandatory for \emph{every} Sparkplug edge node --- \textbf{Node
Control/Rebirth}. It exists to solve a problem that is otherwise unsolvable: a host that
connects \emph{after} a node has already birthed receives data messages it has no birth
certificate to interpret. It knows a value arrived; it does not know the metric's name,
type or unit. Births are published with retain false, so the broker has kept no copy.

The handshake:

\begin{enumerate}
  \item The node's NBIRTH declares a metric named \code{Node Control/Rebirth}, of type
        \code{Boolean}, with the value \code{false}. If the node uses aliases, this metric
        must \emph{not} have one --- a host that has never seen the birth could not know
        the alias.
  \item The host publishes an NCMD containing that metric with the value \code{true}.
  \item The node immediately \textbf{stops sending data}, publishes a complete birth
        sequence (NBIRTH, then every DBIRTH), and only then resumes.
  \item The rebirth reuses the \textbf{same} \code{bdSeq}: no new MQTT session was
        established, so the session number has not changed.
\end{enumerate}

A host also sends a rebirth request when it detects a sequence-number gap that does not
resolve within its configured \emph{reorder timeout} --- the mechanism from
section~\ref{sec:sp-session} feeding directly into this one.

\begin{warning}
An edge node that does not implement \code{Node Control/Rebirth} cannot be repaired by the
host. This was a real gap in \prog until recently; it is now closed --- \prog declares the
metric in every birth certificate and answers a request with a complete birth sequence.
What remains open is the \emph{other} half of the repair: \prog reads no \code{STATE}
topic, so it never learns a host has returned and will not volunteer a birth. The repair
works when the host asks, and only then. See section~\ref{sec:sp-where-we-sit} and
chapter~\ref{ch:contract}.
\end{warning}

\section{The Primary Host and STATE}
\label{sec:sp-primary-host}

A host application announces its own availability, using the same birth/death machinery as
an edge node but on its own topic and in a different format.

\begin{center}
\begin{tabular}{@{}l p{9.6cm}@{}}
\toprule
\textbf{Topic}   & \code{spBv1.0/STATE/sparkplug\_host\_id} \\
\textbf{Payload} & JSON UTF-8 --- not protobuf --- with exactly two keys:
                   \code{online} (a boolean) and \code{timestamp} (UTC milliseconds). \\
\textbf{QoS}     & 1 \\
\textbf{Retain}  & \textbf{true} --- the one place Sparkplug requires retain, so that a
                   node connecting later immediately learns the host's state. \\
\bottomrule
\end{tabular}
\end{center}

The host registers the \code{online: false} version as its \emph{will} in its CONNECT
packet, then --- after subscribing to its own STATE topic --- publishes the
\code{online: true} version. Both carry the \textbf{same} timestamp: the one built at
connect time. So the timestamp identifies the \emph{session}; it is not the moment the
message was published.

\begin{tip}
That equality is a usable diagnostic. If you observe an \code{online: false} whose
timestamp equals the preceding \code{online: true}, it is the \emph{will} --- the broker
published it. If the timestamp is later, the host published it deliberately on its way
out. The two are otherwise indistinguishable on the wire.
\end{tip}

\subsection{What an edge node does with it}

An edge node may be \emph{configured} to treat one host as primary. Nothing forces it to;
Sparkplug explicitly allows a node to have no primary host. But if it is so configured:

\begin{itemize}
  \item it must \textbf{wait} for an \code{online: true} STATE from that host before
        publishing its births;
  \item it must match the host id as the last topic token against its configured value;
  \item it must check that \code{online} really is boolean \code{true};
  \item it must check the timestamp is not older than the last STATE it accepted --- an
        anti-replay rule, because a delayed will from an old session can arrive after a
        new session has started;
  \item and on a valid \code{online: false} it must publish its NDEATH, disconnect, and
        begin session establishment again.
\end{itemize}

\subsection{Why this exists: stranding}

The mechanism looks heavy for a single broker, and it is --- it was designed for
redundancy. With two or more brokers, a node can be perfectly connected to broker~A while
the host has lost \emph{its} connection to broker~A and is happily using broker~B. The
node is publishing into a void and has no way to know. This is called \textbf{stranding},
and STATE is the cure: the node sees the host go offline \emph{on this broker} and walks
to the next one.

\begin{center}
\begin{tikzpicture}[node distance=8mm]
  \node[actor] (n) {Edge Node};
  \node[brokeract, right=20mm of n, yshift=10mm]  (b1) {Broker A};
  \node[brokeract, right=20mm of n, yshift=-12mm] (b2) {Broker B};
  \node[actor, right=20mm of b1, yshift=-11mm] (h) {Host};
  \draw[msg] (n) -- node[above, font=\sffamily\scriptsize, sloped]{connected} (b1);
  \draw[msgdead] (b1) -- node[above, font=\sffamily\scriptsize, sloped]{\emph{lost}} (h);
  \draw[msgback] (b2) -- node[below, font=\sffamily\scriptsize, sloped]{connected} (h);
  \node[annot, below=11mm of n, align=center]
       {without STATE the node keeps publishing to Broker~A\\ and nobody is listening --- \emph{stranded}};
\end{tikzpicture}
\end{center}

\begin{note}
\textbf{On this installation there is one broker, so stranding cannot arise} --- but the
\emph{other} half of the mechanism still bites, because a host restart is a state
transition whether or not there is redundancy. That is measured, not assumed; see
\code{docs/primary-host-state-observation.md} and section~\ref{sec:sp-where-we-sit}.
\end{note}

\section{Conformance profiles}

The specification defines four target roles, and an implementation normally claims exactly
one:

\begin{center}
\begin{tabular}{@{}>{\raggedright\arraybackslash}p{3.7cm} p{8.4cm}@{}}
\toprule
\textbf{Profile} & \textbf{What it must do} \\
\midrule
Sparkplug Edge Node & Manage the session and namespace described above.
                      \textbf{\prog's profile.} \\
Sparkplug Host Application & Consume edge-node messages, and publish its own \code{STATE}
                      birth and death certificates. \textbf{Ignition's profile.} \\
Sparkplug Compliant MQTT Server & Support publish/subscribe at QoS~0 and~1, will messages
                      in full (including retain and QoS~1), and the retain flag in full.
                      Any complete MQTT 3.1.1 or 5.0 broker qualifies. \\
Sparkplug Aware MQTT Server & All of the above, \emph{plus} storing the most recent NBIRTH
                      and DBIRTH of every node and republishing them as retained messages
                      on \code{\$sparkplug} topics. \\
\bottomrule
\end{tabular}
\end{center}

\begin{note}
The \emph{Aware} profile is the standard's own answer to the late-joining consumer: births
that are otherwise unretained become collectable. Mosquitto is Compliant but not Aware, so
on this installation that safety net is absent and the rebirth command
(section~\ref{sec:sp-rebirth}) is the only repair mechanism available. \prog does implement
that command, so the mechanism exists --- but it has to be exercised by the host, which the
Aware profile would not have required.
\end{note}

\section{Where \prog sits}
\label{sec:sp-where-we-sit}

Everything above describes the protocol. This table maps it onto what \prog actually does,
so that nothing in this chapter can be mistaken for a claim about the product.
Chapter~\ref{ch:contract} is the authority, and \code{docs/sparkplug-conformance.md} holds
the clause-by-clause audit --- 303 requirements, each with a verdict and its evidence.

\begin{center}
\begin{tabular}{@{}>{\raggedright\arraybackslash}p{4.3cm} l p{6.6cm}@{}}
\toprule
\textbf{Mechanism} & \textbf{Status} & \textbf{Notes} \\
\midrule
Topic namespace, device level & implemented & One node, one device per meter, keyed on the
                                meter's serial number. \\
NBIRTH / DBIRTH / DDATA / NDEATH & implemented & \code{NDATA} and \code{DDEATH} are not
                                emitted --- see chapter~\ref{ch:contract}. \\
Will registered inside CONNECT & implemented & At \textbf{QoS~0}, where the specification
                                requires QoS~1. A known deviation. \\
Explicit death on a planned stop & implemented & Both certificates fire; a consumer sees
                                two NDEATHs per graceful stop. \\
\code{seq} numbering and wrap & implemented & \\
\code{bdSeq} pairing & implemented & Incremented \textbf{per CONNECT}, as the
                                specification requires, with the will re-registered to match.
                                See chapter~\ref{ch:contract}. \\
Metric properties (quality, units) & implemented & Quality codes \textbf{deviate
                                deliberately} from the specification's enumeration, because
                                the specified codes were measured to display as
                                \emph{good} on Ignition. See chapter~\ref{ch:contract}. \\
Report by exception & \textbf{absent} & Publishes on a fixed interval. Recorded as a
                                deviation. It \emph{was} blocked by the missing rebirth,
                                because periodic publishing substituted for it; that
                                blocker has lifted, and the deviation is now a decision
                                still to be taken rather than one that cannot be. \\
\code{NCMD} subscription & implemented & \code{spBv1.0/<group>/NCMD/<node>} at QoS~1, on
                                every connect, before the birth. Commands other than the
                                one below are logged and discarded --- see
                                chapter~\ref{ch:contract}. \\
\code{Node Control/Rebirth} & implemented & Declared in every \code{NBIRTH}, and answered
                                with a complete birth sequence under an unchanged
                                \code{bdSeq}. Recognised on the name \emph{and} the boolean
                                value \code{true}; any other value is refused and logged
                                with what arrived. \\
\code{DCMD} (writing to outputs) & \textbf{absent} & Nothing is built. Note that two of the
                                four meters \emph{do} expose a relay --- the earlier claim
                                that none did was wrong. Issue~\#38. \\
Primary Host / \code{STATE} & \textbf{declined} & Ruled out deliberately in ADR 0018, not
                                merely unbuilt. See chapter~\ref{ch:contract}. \\
Aliases, Templates, DataSets & \textbf{absent} & Out of scope; metric names are published in
                                full. \\
Multiple brokers & \textbf{absent} & One broker is configured; server-walking cannot arise. \\
\bottomrule
\end{tabular}
\end{center}

\begin{warning}
\textbf{One of the two gaps is now closed, and the argument has to be re-run rather than
re-worded.}

The situation this box described was a compound one. \prog's NBIRTH is published once, at
QoS~0 and not retained, so the broker keeps no copy for a returning consumer. It reads no
\code{STATE} topic, so it never learns that a host has restarted and never re-births on
that account. Its own broker session is untouched by a host restarting, so nothing prompts
it to reconnect. And --- the leg that has changed --- the protocol's remaining repair, a
rebirth request, reached \prog and was ignored.

\textbf{It is no longer ignored.} \prog declares \code{Node Control/Rebirth} in every birth
certificate and answers a conformant request with a complete birth sequence. So there is now
a host-initiated path back, and Ignition's MQTT Engine takes it on its own: receiving data
from a node whose birth it never saw is exactly the condition under which it issues a
rebirth request.

\textbf{What still does not exist is a \prog-initiated path.} The bridge cannot notice that
a host has returned, because it reads no \code{STATE}. If a host does not ask, nothing
happens. The remaining exposure is therefore narrower and differently shaped: not \emph{no
repair exists}, but \emph{the repair depends on the consumer asking for it}. A
\emph{Sparkplug Aware} broker would have removed that dependency by replaying the stored
births; Mosquitto is not one.

\textbf{Any reconnect on \prog's own side also restores the definitions}, since a fresh
NBIRTH is published on every successful MQTT connection --- a broker restart, a network
interruption or a keep-alive expiry all repair the consumer's view. What no event on the
\emph{host} side can do is prompt that reconnect. (An earlier edition of this paragraph
said the definitions could not be restored \emph{until \prog itself is restarted}. That was
false. The same sentence was corrected in chapter~\ref{ch:contract} first, and this copy of
it survived the correction by a chapter.)

\textbf{The consequence for planning:} the rebirth handling was ranked ahead of the
\code{STATE} work because it was the cheaper and more load-bearing of the two. That
reasoning has now been spent, and the \code{STATE} decision has to be weighed on its own
evidence rather than inheriting a ranking that was argued from a gap which no longer
exists. The measurement it will be weighed against was taken on the live installation
before anything was designed.
\end{warning}

\section{Reading the specification yourself}
\label{sec:sp-reading-the-norm}

The Sparkplug~B specification is vendored in this repository, at

\begin{center}\code{docs/spec/sparkplug-b-3.0.0/}\end{center}

\noindent release tag v3.0.0, EPL-2.0. Vendoring it means it can be grepped and cited
without a network connection, and that it cannot drift underneath the audit that
references it.

Each testable requirement carries an identifier of the form

\begin{center}\code{tck-id-payloads-ndeath-will-message-qos}\end{center}

\noindent Citing that identifier, rather than paraphrasing the prose around it, is the
convention used throughout this project's documentation and issue tracker.

\begin{warning}
\textbf{This project's standing rule: read the norm, not a summary of it.} Both of the
worst defects found so far came from trusting a secondary source --- a vendor
documentation page and a plausible recollection --- rather than the specification text.
One published quality codes that a real host displayed as \emph{good}, which is precisely
the failure this bridge exists to prevent; the other rejected a legal design as a
``specification violation'' on the strength of a requirement the specification does not
contain.
\end{warning}
